% =========================================================================
% REFERÊNCIAS
% =========================================================================
\clearpage
\begingroup
\singlespacing
\chapter*{Referências}
\addcontentsline{toc}{chapter}{REFERÊNCIAS}

\noindent
ABNT -- Associação Brasileira de Normas Técnicas. \textbf{NBR 14724}:
Informação e documentação. Trabalhos Acadêmicos -- Apresentação. Rio de
Janeiro: ABNT, 2002.

\vspace{0.5\baselineskip}

\noindent
BOYER, C. B.; UTA, C. M. \textbf{História da Matemática} [Trad. Helena
Castro]. 3.~ed. São Paulo: Blucher, 2012.

\vspace{0.5\baselineskip}

\noindent
D'AMBRÓSIO, U. \textbf{Educação Matemática:} da teoria à prática. 23.~ed.
Campinas: Papirus, 2012.

\vspace{0.5\baselineskip}

\noindent
KUBO, O.; BOTOMÉ, S. \textbf{Ensino e aprendizagem:} uma interação entre
dois processos comportamentais. Interação, v.~5, p.~123--32, 2001.

\vspace{0.5\baselineskip}

\noindent
HART-DAVIS, A. \textbf{O Livro da Ciência.} 2.~ed. São Paulo: Globo, 2016.

\vspace{0.5\baselineskip}

\noindent
PILETTI, C. \textbf{Didática geral.} São Paulo: Ática, 1995.

\vspace{0.5\baselineskip}

\noindent
RIBEIRO, J. L. P. Áreas e Proporções nas Superquadras de Brasília Usando o
Google Maps. \textbf{Revista do Professor de Matemática}. Rio de Janeiro,
n.~92, p.~12--15, jan-abr. 2017.

\vspace{0.5\baselineskip}

\noindent
SEVERINO, A. J. \textbf{Metodologia do trabalho científico.} 22.~ed. rev. e
ampl. São Paulo: Cortez, 2002.
\endgroup
